\section{Discussion, falsifiable claims and research agenda}\label{sec:discussion}

\subsection{What has been proved}
The objective \eqref{eq:objective} remains the classical PLS objective.
Theorem~\ref{thm:odds} characterizes exactly which shrinkage
weights belong to a single quadratic-penalty family;
Theorem~\ref{thm:root} proves a unique scalar root for
every prescribed interior $s$ and provides a finite bracket;
Theorem~\ref{thm:reconstruction} proves exact equivalence
between penalty-based and direct-weight fits, including
fixed nonzero means of differences and both endpoints.
Theorem~\ref{thm:sensitivity} quantifies sensitivity and
separates weight stability from scalar-root conditioning.
Theorem~\ref{thm:plugin} supplies the full Jacobian trace
when the difference mean is computed from the observations.
The results are algebraic consequences of an established
linear smoother. These are algebraic results for an established estimator;
their presentation does not establish historical priority.

\subsection{Why using weights can nevertheless be useful}
For an interior target $s$ and a known positive spectrum,
the weight representation computes all shrinkage factors
from a finite scalar $x$. In exact arithmetic,
$\lambda=\exp(x)/\delta_*$ from \eqref{eq:implicit};
the direct-weight fitting path has no need to evaluate
that expression. It may therefore be possible to
reconstruct a trend in a regime where $\lambda$ itself
cannot be represented as a finite floating-point
number, provided the singular values and the weight
root remain numerically resolvable. This is a possible advantage of the representation,
not a general proof of numerical superiority. If the smallest singular
values have already been corrupted or lost to underflow,
the method cannot recover that information.

Caching $w(s)$ for repeated series with identical $(n,d,s)$
also avoids repeated inversion. A well-optimized penalty-based implementation can cache
the same weights. The new parameterization therefore has
no automatic asymptotic advantage. The correct comparison distinguishes
(i) initial SVD cost, (ii) scalar inversion, (iii) repeated
full-vector reconstruction, (iv) memory, and (v) failed
or inaccurate candidate fits. Prior structured
Whittaker--Henderson algorithms must remain in the
comparison \citep{weinert2007,biessy2026}.

\subsection{Proposed Experiment006: a falsifiable validation}
A later \emph{separately numbered} experiment should
compare three implementations on identical inputs:
(A) stable log-$\lambda$ inversion followed by
shift-free spectral fitting; (B) the present direct
log-odds weight fit; and (C) the direct-weight path
with cached weights and repeated observations.
Use $d=1,\ldots,6$, multiple $n$ values,
fixed $\mu=0$ and fixed nonzero $\mu$, as well
as the separately labeled observed-difference
plug-in. Cover an interior $s$ grid and attainable
values near zero and one.

At every identical $(n,d,s,\mu,y)$,
measure total and phase-separated time, root
evaluations, allocation and peak memory,
attained smoothness, method failure, numerical
relative error against an independent
high-precision \emph{same-estimator} reference,
and statistical error against known latent trends
when synthetic truth is available. Include the
same scalar precision tolerance and a
cache-equivalent penalty-based baseline.
Block independent replicates by simulated
series and machine run; do not treat hundreds
of $s$-grid points as independent samples.
Failure counts should remain in the raw data.
A direct-weight advantage would be supported by
matched measurements and high-precision accuracy
at the same requested $s$, not by an isolated
microbenchmark or a faster result from a
different candidate grid.

Two especially testable predictions are:
(i) a direct-weight path may avoid forming an
overflowing $\lambda$ in some high-order,
high-smoothness cases where the log-singular
spectrum remains usable; and (ii) caching weights
will be most beneficial when one $(n,d,s)$
configuration is reused across many observations.
Neither prediction is claimed as verified here.
An equally optimized log-$\lambda$ and weight-cache
baseline might match both outcomes.

\subsection{Limits and scope for a standalone paper}
Our results concern unweighted, equally indexed
finite differences and quadratic penalties.
They do not establish a theory of irregular
observation spacings, time-varying smoothness,
$\ell_1$ trend filtering, generalized likelihoods,
or a new forecast-selection criterion.
The smoothness coordinate is a normalized
\emph{conditional trace index}; when $\mu$
is estimated from data the full fitted-value
Jacobian has the additional term in
\eqref{eq:edf-plug}. If $s$ is selected from
the data, its adaptive effect requires
further analysis. A bounded parameter interval
does not imply a globally well-conditioned
inverse or a globally optimal finite search.

\section{Conclusion}
A PLS trend can be reconstructed by solving
one monotone scalar equation for a coupled vector
of spectral weights and applying a shift-free SVD
formula, without explicitly calculating the
penalty coefficient in the fitted-trend path.
The estimator, its shrinkage family, and the
implicit scalar inversion remain those of the
classical penalized problem. The exact
endpoint projection, sensitivity bound and
plug-in degrees-of-freedom correction help
clarify what prescribing $s\in[0,1]$
mathematically controls. The proposed paired precision--runtime study is needed
to determine whether direct reconstruction improves
accuracy or execution time in practice.

\section*{Software and evidence status}
The implementation is maintained at
\url{https://github.com/heri-espino/trendestimation}.
Relevant sources are the research module
\texttt{smoothness\_weights.py} and corresponding
\texttt{test\_spectral\_smoothness\_weights.py}.
This document is a theory-first research
draft. Its theorems have proofs, but no new
runtime or statistical Monte Carlo measurements
were conducted for the manuscript.
Authorship, journal target and statistical
claims are provisional pending independent
coauthor review.
